\documentclass[12pt]{article}
\usepackage{styles/cwilliams-standard}

\setclass{IPD}
\settitle{Smelly Code}
\setauthors{Mihir Patel, Arjun Sivakumar, Christopher Williams}
\setheaderauthors{Patel, Sivakumar, Williams}

\begin{document}

\maketitlepage

\begin{aside}{Layout}
    The layout of this document goes by section, with the specific ``smell'' listed on 
    the table of contents page for easier navigation. The group member that did what smell 
    is listed in each section.
\end{aside}

\tableofcontents
\newpage

\section{Bloaters}

\textit{Mihir}

\subsection{Large Class}

\begin{description}
    \item[Location in Code]
        File: \texttt{Python/SmellAnnotated/Shop.py} \\
        Class: Shop (overall class; spans multiple methods and fields)

    \item[What is happening here?]
        The Shop class acts like a ``do everything'' class: it creates pizzas, receives orders,
        stores customer contact info, sends notifications, applies discounts and loyalty points,
        and handles complaints---all inside one class.

    \item[Concrete Evidence]
        The class contains many distinct responsibilities across its methods:
        \begin{itemize}
            \item Pizza creation: \texttt{create\_pizza} (lines 23--26)
            \item Order intake: \texttt{receive\_order} (lines 38--43)
            \item Customer contact management: \texttt{update\_contact\_info}, \texttt{update\_name},
                  \texttt{update\_address}, \texttt{update\_phone\_number}, \texttt{update\_email} (lines 45--72)
            \item Customer notifications: \texttt{notify\_for\_promotion}, \texttt{notify\_for\_discount}, \\
                  \texttt{notify\_for\_new\_arrivals} (lines 74--86)
            \item Benefits: \texttt{apply\_discount}, \texttt{apply\_loyalty\_points} (lines 88--96)
            \item Complaint handling: \texttt{handle\_complaint} (lines 98--109)
        \end{itemize}
        Additionally, customer contact fields (\texttt{first\_name}, \texttt{last\_name}, \texttt{address}, \\
        \texttt{phone\_number}, \texttt{email}) and temporary order fields (\texttt{temp\_discount\_code},
        \texttt{temp\_order\_note}) are stored directly on Shop. The class also imports and directly
        coordinates Cashier, Chef, and multiple pizza subclasses (CheesePizza, VeggiePizza,
        TunaPizza, PepperoniPizza) via \texttt{self.chef}, \texttt{self.casher}, and \texttt{self.pizza}.
        Key evidence checklist items: has many dependencies/imports; very large class; exposes or
        relies on internals of another module.

    \item[Why is this risky?]
        A large class tends to have low cohesion---unrelated concerns are mixed together---making it
        harder to test and making changes more likely to cause side effects in seemingly unrelated behavior.

    \item[Estimated Change Radius]
        Around 3--6 files for a typical customer/order flow change: \texttt{Shop.py} plus the classes
        it coordinates (\texttt{Cashier.py}, \texttt{Chef.py}, \texttt{Pizza.py} or pizza subclasses).

    \item[Impact]
        More bugs during maintenance because developers must understand many features at once, and
        slower changes since updates to one feature (like discounts or complaints) are entangled
        with the same central class.

    \item[Confidence Level]
        \textbf{High.} The class visibly contains many distinct responsibilities and directly
        orchestrates multiple other components (Cashier, Chef, multiple pizza subclasses).

    \item[Suggested Refactor]
        Use Extract Class / responsibility split: move customer contact fields and update methods
        into a \texttt{CustomerProfile} class, move notification methods into a \\
        \texttt{Notifier}/\texttt{MarketingService}, and move discount/loyalty logic \\ 
        into a \texttt{DiscountPolicy}/\texttt{LoyaltyProgram}, leaving Shop to coordinate only
        high-level order flow.
\end{description}

\section{Couplers}

\textit{Arjun}

\subsection{Feature Envy}

\begin{description}
    \item[Location in Code]
        File: \texttt{shop.py} \\
        Methods and lines:
        \begin{itemize}
            \item \texttt{create\_pizza()}, \texttt{receive\_order()} --- lines 27--42
            \item \texttt{handle\_complaint()} --- line 79
            \item \texttt{chain\_of\_methods()} --- line 98
            \item \texttt{access\_internal\_details()} --- line 113
        \end{itemize}

    \item[What is happening here?]
        The Shop class uses features and internal data of multiple other classes (Cashier, Chef,
        and Pizza) and relies on their internal functions and data more than its own, making Shop
        deeply coupled with those classes.

    \item[Concrete Evidence]
        \textbf{Evidence 1:} \texttt{create\_pizza} and \texttt{receive\_order} \\
         (lines 27--42 of \texttt{shop.py}) instantiate Pizza subclasses directly (\texttt{CheesePizza},
        \texttt{VeggiePizza}, \texttt{TunaPizza}, \texttt{PepperoniPizza}) based on assumptions
        about the Pizza inheritance hierarchy, and then delegate to \texttt{self.casher.take\_order}.
        The logic properly belongs inside the Pizza or Cashier classes.

        \textbf{Evidence 2:} \texttt{access\_internal\_details} (line 113) traverses two levels of
        object internals---\texttt{self.casher.chef.busy}---to read a private field deep inside
        Chef via Cashier. This is a textbook Law of Demeter violation. Key evidence checklist items:
        touches many files for one change; has many dependencies/imports; very large class; breaks
        tests or features far away from change site; exposes or relies on internals of another module.

    \item[Why is this risky?]
        Deep coupling means \texttt{shop.py}, \texttt{cashier.py}, \texttt{chef.py}, and
        \texttt{pizza.py} cannot independently evolve. A change to an internal field or method in
        \texttt{chef.py} can cause \texttt{shop.py} to crash or behave unexpectedly, making the
        system fragile and hard to test in isolation.

    \item[Estimated Change Radius]
        Up to 4 files for a typical change: \texttt{shop.py} plus whichever of \\
         \texttt{cashier.py}, \texttt{pizza.py}, and \texttt{chef.py} the modified function depends on.

    \item[Impact]
        Increased defect rates and maintenance costs---unexpected behavior propagates across class
        boundaries whenever internal data or behavior of a dependency changes.

    \item[Confidence Level]
        \textbf{Medium.} The dependencies are not deeply convoluted and a few targeted fixes could
        make Shop more independent, but the Feature Envy smell is clearly present.

    \item[Suggested Refactor]
        Remove direct reliance on other classes' internals. At minimum, wrap dependencies in
        try/catch blocks to prevent cascading crashes. Ideally, extract an \texttt{OrderManager}
        class to consolidate cross-class coordination, and push pizza-creation logic into a factory
        method on the Pizza hierarchy so Shop no longer needs to know about subclass details.
\end{description}

\section{Change Preventers}

\textit{Chris}

\subsection{Shotgun Surgery}

\begin{description}
    \item[Location in Code]
        Files: \texttt{Cashier.py}, \texttt{Chef.py}, \texttt{Customer.py}, \texttt{Pizza.py}, \texttt{Shop.py} \\
        Classes: Cashier, Chef, Customer, Pizza, Shop \\
        Methods: \texttt{update\_contact\_info}, \texttt{update\_name}, \texttt{update\_address}, \\
        \texttt{update\_phone\_number}, \texttt{update\_email}

    \item[What is happening here?]
        A group of methods that all address/interact with the same information.

    \item[Concrete Evidence]
        In all 5 files and classes listed above, each class has the listed methods that all
        interact with information (name, address, phone number, email) and split the updates
        to them among \texttt{update\_contact\_info} as well as individual methods that address each
        individual piece of information. 

        This smell:
        \begin{itemize}
            \item Touches many files for one change (large change radius)
            \item Duplicates logic found elsewhere
            \item Breaks tests or features far away from change site
        \end{itemize}

    \item[Why is this risky?]
        Making one change requires making many smaller changes to many classes, and it
        splits a single responsibility among many different classes and methods.

    \item[Estimated Change Radius]
        All 5 classes and methods---around 25 spots in the code in total.

    \item[Impact]
        Worse organization, more code duplication, and harder maintenance.

    \item[Confidence Level]
        \textbf{High.} They all have the same code reused doing the same function, split among
        different classes and methods.

    \item[Suggested Refactor]
        {Extract class}. Move responsibility to a superclass that all methods could extract/pull
        from, or create an interface that all classes can use.
\end{description}

\section{Dispensables}

\textit{Arjun}

\subsection{Duplicate Code}

\begin{description}
    \item[Location in Code]
        File: \texttt{chef.py}, lines 46--98 \\
        Duplicate methods present in \texttt{chef.py} that also exist in \texttt{cashier.py} and \texttt{pizza.py}:
        \begin{itemize}
            \item \texttt{update\_contact\_info}, \texttt{update\_name}, \\ 
            \texttt{update\_address}, \texttt{update\_phone\_number}, \texttt{update\_email}
            \item \texttt{notify\_for\_promotion}, \texttt{notify\_for\_new\_arrivals}
            \item \texttt{apply\_discount}, \texttt{apply\_loyalty\_points}
            \item \texttt{handle\_complaint}, \texttt{ask\_for\_receipt}
            \item \texttt{another\_unused\_method}, \texttt{yet\_another\_unused\_method}
        \end{itemize}

    \item[What is happening here?]
        \texttt{chef.py} contains multiple functions (lines 46--98) that are exact \\
        copies of code found in \texttt{pizza.py} and \texttt{cashier.py}. These functions make sense in
        \texttt{cashier.py} (e.g., updating contact info or notifying customers), but have no
        business belonging to the Chef class.

    \item[Concrete Evidence]
        \textbf{Evidence 1:} \texttt{update\_contact\_info} at line 46 of \texttt{chef.py} and
        line 84 of \texttt{cashier.py} are character-for-character identical---same signature,
        same field assignments (\texttt{self.first\_name}, \texttt{self.last\_name},
        \texttt{self.address}, \texttt{self.phone\_number}, \texttt{self.email}).

        \textbf{Evidence 2:} \texttt{ask\_for\_receipt} at line 91 of \texttt{chef.py} and
        line 38 of \texttt{cashier.py} are identical: both simply print
        \texttt{"Customer is asking for a receipt."} Key evidence checklist items: very large class
        or very long method; duplicates logic found elsewhere; bypasses architectural layers.

    \item[Why is this risky?]
        The repeated code reduces modularity and makes the codebase less testable and cohesive,
        since the same edge cases must be covered for each copy. It also violates the DRY
        (Don't Repeat Yourself) and YAGNI (You Aren't Gonna Need It) principles, reducing
        maintainability and testability.

    \item[Estimated Change Radius]
        Depends on which functions are changed. Modifying duplicate-only code is locally contained,
        but the Chef class also has core functions used by other classes (e.g., \texttt{Cashier}
        calls \texttt{self.chef.bake\_pizza} inside \texttt{Cashier.take\_order}), meaning some
        changes could ripple outward.

    \item[Impact]
        Testing and maintenance become more difficult as the class grows unnecessarily complex and
        error-prone. Bug fixes or tests for a given function must be applied to every copy, multiplying
        time and effort.

    \item[Confidence Level]
        \textbf{High.} The duplicate code is very clearly a repetitive chunk that serves no purpose
        and makes no sense under the Chef class.

    \item[Suggested Refactor]
        Remove the copies of the misplaced functions from \texttt{chef.py}. Responsibility for
        customer-facing concerns (managing contact details, notifications, discounts) should be
        moved entirely to \texttt{cashier.py} or a dedicated service class, not duplicated across
        unrelated classes.
\end{description}

\section{Object-Oriented Abusers}

\textit{Chris}

\subsection{Refused Bequest}

\begin{description}
    \item[Location in Code]
        Files: \texttt{Customer.py}, \texttt{Cashier.py}, \texttt{Chef.py}, \texttt{Pizza.py} \\
        Classes: Customer, Cashier, Chef, Pizza \\
        Method: \texttt{refused\_bequest(self)}

    \item[What is happening here?]
        A method from a superclass that should be overridden in subclasses, but is left empty.

    \item[Concrete Evidence]
        All 4 files listed above contain a \texttt{refused\_bequest} method that does nothing.

        This smell:
        \begin{itemize}
            \item Bypasses architectural layers
            \item Exposes or relies on internals of another module
        \end{itemize}

    \item[Why is this risky?]
        This method is intended to be overridden but is not used in the base class/superclass,
        indicating the subclasses are inheriting behavior they have no relation to.

    \item[Estimated Change Radius]
        Every class that contains the \texttt{refused\_bequest} method.

    \item[Impact]
        Confusion and delays---developers must wonder why a class inherits a method from a
        superclass that it is not related to and does not use.

    \item[Confidence Level]
        \textbf{Medium.} The name \texttt{refused\_bequest} and that method going unused make it
        identifiable, though the intent behind the inheritance is ambiguous.

    \item[Suggested Refactor]
        Either replace inheritance with delegation, or extract the method out of the superclass
        into a separate new superclass, and have both the subclass and original superclass
        inherit from that new superclass instead.
\end{description}

\end{document}